\section{Discussion}
\label{sec:discussion}

\subsection{What the Evidence Supports}

The evaluation addresses the three objectives through distinct but interacting components: maintaining record-level detection under legitimate change, interpreting batch changes, and allocating validation work. The results support conditional benefits rather than uniform superiority of the complete adaptive configuration.

\paragraph{The gated drift signal carries the batch labels.} Removing the gate reduces macro-F1 by roughly 0.23--0.31 across all datasets. Filtering records that fail basic checks keeps some error-induced distribution movement out of the drift estimate, allowing failure evidence and distribution evidence to play different roles. On Chicago, fixed or EWMA failure thresholds behind the same gate produce identical predictions to adaptive thresholds. On Retail, their point estimates are higher and the paired analysis establishes no adaptive advantage. Adaptation adds most clearly on BTS, where fixed bounds otherwise label 89\% of pure-drift batches as degraded. Thus the strong four-way scores should be attributed to the combination of signals and their calibration, with adaptation's incremental contribution assessed separately. The gate is still heuristic: excluded records may include valid novelty, and undetected errors may remain in the retained data.

\paragraph{Adaptation pays off where drift moves valid values across fixed bounds.} In BTS, a changing delay or taxi-time distribution increasingly crosses fixed limits; adjusting location and scale removes much of this additional false alarming. In Retail, broad fixed ranges already accommodate the tested multiplicative change. Adaptation then trades a small FPR increase for higher recall, which does not satisfy an FPR-reduction hypothesis. Chicago is different again: its numerical columns retain physical bounds, and the improvement comes from category recurrence. These cases show why policy performance depends on both the data distribution and the monitored domain. In practice, a historical replay can establish whether fixed limits produce growing false alarms under expected change before numerical adaptation is selected.

\paragraph{Adaptation absorbs persistent errors.} Repeated observations are evidence of frequency, not independent evidence of validity. The tested numerical policy loses up to 29.06 points of recall under persistent taxi-out contamination. The categorical policy admits recurring invalid codes, reducing statistical semantic recall from approximately 0.51--1.00 under fixed rules to 0.07--0.13 under adaptive rules. All profiles enter history, so the drift gate does not protect numerical learning from suspect batches. Fixed domain constraints and guarded category admission can retain distinctions that recurrence and broad ranges lose. Their evidence must nevertheless remain configuration-specific: fixed+LLM rules attain the highest injected semantic recall, while the per-value guard is evaluated separately on only 24 distinct values. No result here establishes a fully integrated guarded monitor's performance.

The connection with poisoning research~\cite{Biggio2012,Steinhardt2017} is an interpretation of the same learning mechanism, not an evaluated adversarial threat model. A persistent invalid value can become statistically normal whether it originates in an accidental mapping error or deliberate manipulation. The present evidence supports limiting what history may authorise, without claiming an evaluated security defence.

\paragraph{Selective execution is limited by monitoring overhead.} Risk selection meets the point-estimate usefulness criterion more often than matched periodic or random schedules, showing that the observed risk signals help allocate some checks. Savings remain modest because basic checks always run and profile and drift work continue on every batch. After reuse, skipping optional checks produces no material end-to-end timing benefit in these workloads. Conversely, the ablation shows that text-pattern rules add scanning work: their main-effect byte difference is 30--38\% of the dataset mean. Selectivity is therefore most relevant when optional checks account for a substantial part of total cost, a condition that must be measured rather than inferred from check counts. Sparse and short incidents expose the opposing risk: skipping a check also withholds the evidence that would raise its priority.

\subsection{Implications for Practice}

Four implications follow directly from the evaluated setting. First, retain separate distribution-change and check-failure evidence. A single unusualness score cannot identify why a batch changed, and an operational label should remain interpretable through the contributing columns and checks. Second, calibrate drift thresholds against the available warm-up variability and verify the resulting stable-batch behaviour. The resampling procedure improved development calibration here, but five historical batches do not justify a general finite-sample guarantee.

Third, distinguish a data contract from a learned expectation. Physical bounds and documented domain membership require independent justification; repeated presence or compatibility with a pattern is insufficient. Historical screening remains useful for rejecting overly restrictive proposals, but relaxing a semantic rule to observed values can also weaken its intended meaning. The dictionary-only model rules demonstrate the value of screening, while their unequal information access prevents a comparison with human expert quality.

Fourth, measure total monitoring cost as well as work avoided. Reuse existing profile statistics before interpreting skipped-check savings as a runtime benefit. Report incident recall and observation staleness alongside the execution rate, because a schedule can save checks while missing the short events that matter operationally. These recommendations apply to the evaluated batch setting; deployment-specific tolerance for missed errors and alarms remains a policy decision.

\subsection{Threats to Validity}

\paragraph{Internal validity.} Development and test pools contain disjoint records and use separate seeds, but derive from the same three datasets. Changes to calibration, retention of known categories and selective-execution gating were made during development and documented before the internal freeze. Three later sensitivity dials are explicitly post-hoc. Manifests identify the tested code and settings, but do not provide independent public preregistration. Comparator implementation remains an author-controlled factor. In particular, the AVH-inspired baseline omits the original optimiser, and default versus calibrated tool configurations should not be interpreted as a tool-wide ranking.

\paragraph{Construct validity.} Injected labels isolate mechanisms but simplify natural error processes. A positive injected drift level defines drift even when the change is operationally small; degradation denotes an increased error rate rather than an independently logged business incident. Rule failures provide observable evidence, not a direct census of all errors. Current-batch adaptation also uses information unavailable to a strictly predictive policy. Record FPR does not fully express batch alert burden. Logical bytes touched approximate implementation work rather than measured cloud I/O, and all timing comes from one seed and one machine.

\paragraph{External validity.} Three domains, two dirty/clean tables and one natural timeline do not represent all pipelines. LLM findings depend on the specific endpoints, prompts and selected violations. The dictionary-only baseline is another model, not a domain expert. Category guarding tests twelve valid and twelve invalid distinct values, repeated across seeds; it does not cover initially present invalid categories, which bypass the per-value check. Tax cell recall is conditional on checkable columns, and the strongest cross-column configuration has low cell precision. Chicago's known procedure change provides one interpretable event, while all 3,276 Facility Type alarms still require independent human adjudication.

\paragraph{Conclusion validity.} Ten seed clusters provide uncertainty about replay variation conditional on the source datasets. Bootstrap intervals can be coarse, and pooled macro-F1 differs from the seed-mean estimand used for its paired comparison. Primary H1a decisions are per scenario, not a study-wide familywise guarantee. The separate secondary Holm reporting correction does not remedy the approximate calibration of the original bootstrap sign-tail $p$-values. Descriptive sensitivity criteria, H3 point-estimate passes and non-significant comparisons should not be read as inferential guarantees. Finally, statistical differences may be operationally minor: Chicago's approximately 0.03 pp FPR reduction is clear but small, whereas the BTS contamination recall losses are substantial.
